% Owner: Prisca.
\section{Threat Model and Research Questions}
\label{sec:threat}

\textbf{Setting.} An LLM reviews a JavaScript web handler and answers whether it is vulnerable, ending with a
verdict line (\vuln{} or \safe{}). Between two reviews the code changes while the vulnerability stays. The change
may be made by the model itself (a maintenance edit), by a tool or developer (a refactoring, a split into
modules), or by an adversary who wants the bug to pass review. The adversary can edit the code under review,
including comments, strings and imports, and in the pushback setting can reply to the model in the conversation.
It cannot change the reviewer's system prompt and has no access to the model's weights or scores.

\textbf{What counts as a miss.} A \emph{miss} is a \safe{} verdict on a program that our verifier confirms is
still exploitable. A program is \emph{safe} if the sink of its vulnerability class cannot be abused: no injected
shell command, no file read outside the base directory, no unescaped input in HTML output. Other weaknesses
(for example server-side request forgery or argument injection) are out of scope for the label
(Section~\ref{sec:limitations}).

\textbf{Research questions.}
\begin{itemize}
\item[\textbf{RQ1}] Do the model's own QA edits (documentation, logging, refactoring, error handling, input
      validation) make it miss a bug it caught? (Part 1)
\item[\textbf{RQ2}] Do semantics-preserving tarpit transformations, alone or stacked, make it miss? (Part 2)
\item[\textbf{RQ3}] Does embedding the bug in a larger, multi-file program make it miss, and what happens to false
      alarms on safe programs? (Part 3)
\item[\textbf{RQ4}] Which deliberate, bug-preserving attacks do make it miss? (Part 4)
\item[\textbf{RQ5}] Does a standard prompt-level defense stop those attacks, and do they transfer to a commercial
      model?
\end{itemize}
